\section{Introduction}
\label{sec:intro}

Truncated or inexact Newton methods solve the Newton system approximately with an
iterative solver rather than forming the Hessian
\citep{dembo1982inexact,steihaug1983cg,nash1984lanczos,nash2000survey}.
Implementations that access curvature only through Hessian-vector products build
on the exact product of \citet{pearlmutter1994hvp}; applications at deep-network
scale are reported by \citet{martens2010hessianfree,martens2011rnn}. Practical
performance depends strongly on two computational-resource decisions: the
damping level and the per-step CG iteration budget.

Damping belongs to the Levenberg--Marquardt family of regularizers
\citep{levenberg1944,marquardt1963} and is related to a trust-region radius
\citep{steihaug1983cg,conn2000trustregion}. In Hessian-free optimization the
damping update rule is reported to drive performance
\citep{martens2010hessianfree,martens2011rnn}. Truncating the CG budget is the
central design variable of inexact Newton methods
\citep{dembo1982inexact,nash2000survey}.

Attempts to adapt these two quantities during optimization rather than fixing
them connect to the learned-optimizer literature
\citep{andrychowicz2016l2l,metz2019pathologies,metz2020effective,bae2022apo}.
For reinforcement learning \citep{schulman2017ppo} to be the right tool for
learning such a controller, observing the state and changing the action in
response must carry value. The distinction between fixing a plan in advance and
reacting to the state during execution is the classical distinction between
open-loop and closed-loop control \citep{bertsekas2017dp}.

Our question is not whether performance improves, but \emph{where the
improvement comes from}.

\begin{quote}
In Hessian-free Newton--CG, does the benefit of state-dependent computational
resource control arise from multi-step planning or from execution-time feedback?
\end{quote}

We split this into two testable propositions.

\begin{description}
  \item[Q1] Does a good action sequence exist?
  \item[Q2] Is it worth revising that sequence with feedback during execution?
\end{description}

The two questions imply different artifacts. If only \textbf{Q1} holds, what is
needed is a predictor that fixes a schedule at the initial state. Only if
\textbf{Q2} holds is a per-step feedback policy justified.

\subsection*{Contributions}

\begin{enumerate}
  \item We operationalize the control of damping and CG effort as a sequential
        decision problem.
  \item We present a controller ladder that separates multi-step planning from
        execution-time feedback.
  \item On held-out instances of ill-conditioned SPD quadratics, we confirm an
        additional value for multi-step planning ($+0.456$~nat, CI $+0.254$ to
        $+0.720$, $p < 0.0001$, $35/40$) and do not observe a practically large
        additional value for feedback ($+0.010$~nat, CI $-0.033$ to $+0.053$).
  \item We document and apply an audit procedure required to make our controller
        comparisons identifiable.
\end{enumerate}

Item~1 is not a new MDP formulation. Item~4 is not a general benchmark
framework; it is a \emph{secondary} contribution, presented as part of the
methods in Section~\ref{sec:eligibility} and revisited in
Section~\ref{sec:audit-why}.

The decision not to train a policy is not listed as a contribution.
Section~\ref{sec:nopolicy} separates the scientific conclusion from the project
go/no-go decision.
